% -------------------------------------------------------
% Secção 2 – Conhecimentos na Fundamentação da Abordagem
% -------------------------------------------------------

% -------------------------------------------------------
\subsection*{2.1 Coordenação de Transações Distribuídas — Modelo X/Open}
\addcontentsline{toc}{subsection}{2.1 Coordenação de Transações Distribuídas — Modelo X/Open}

Quando múltiplos serviços com estado participam na mesma unidade de trabalho, a questão central
é como garantir que todos confirmam ou todos revertem as suas alterações --- mesmo em presença de
falhas parciais de rede ou de nós. A abordagem de \textbf{coordenação centralizada}, em que um
componente dedicado (o TM) dirige o protocolo de commit, oferece semânticas fortes e um modelo
de falha bem definido, à custa de tornar o TM um potencial ponto único de falha. A
\textbf{coordenação distribuída}, em que cada participante toma decisões autónomas, evita este
ponto único mas exige consenso entre os participantes, aumentando a complexidade e a latência
do protocolo. O modelo X/Open DTP adota a abordagem centralizada, separando claramente as
responsabilidades de cada componente através de interfaces padronizadas.

O modelo \textit{X/Open Distributed Transaction Processing} (DTP) define uma arquitetura
padronizada para a coordenação de transações em sistemas distribuídos, assente em três
componentes com responsabilidades bem delimitadas:

\begin{itemize}
  \item \textbf{Application Program (AP)} --- inicia e termina transações via interface TX;
        no sistema implementado, corresponde ao \texttt{Service.Client} e ao \texttt{Service.Validator}.
  \item \textbf{Transaction Manager (TM)} --- coordena o protocolo 2PC entre os RMs participantes;
        implementado como \texttt{Service.TM}.
  \item \textbf{Resource Manager (RM)} --- gere um recurso transacional e expõe a interface XA;
        implementado como \texttt{Service.Vector}.
\end{itemize}

O protocolo \textit{Two-Phase Commit} (2PC) é o mecanismo pelo qual o TM garante atomicidade
distribuída. Na \textbf{Fase 1 — Prepare}, o TM envia \texttt{xa\_prepare()} a cada RM participante;
cada RM persiste em armazenamento estável toda a informação necessária para commit ou rollback e
responde afirmativamente (\texttt{XA\_OK}) ou negativamente (\texttt{XA\_RB*}), sendo que uma
resposta negativa veta imediatamente a transação global. Na \textbf{Fase 2 — Commit/Rollback},
se todos os RMs votaram afirmativamente, o TM regista a decisão de commit em log persistente e
envia \texttt{xa\_commit()} a todos; se qualquer RM vetou, o TM envia \texttt{xa\_rollback()} a
todos os restantes. A política de \textit{presumed rollback} — na ausência de registo para uma
transação, assume-se que foi revertida — reduz o número de escritas em armazenamento estável no
caminho mais comum.

O TM constitui um potencial ponto único de falha: se falhar após completar a Fase~1 mas antes de
concluir a Fase~2, os RMs participantes ficam bloqueados num estado \textit{in-doubt}, incapazes
de libertar recursos. A mitigação é um \textbf{log de transações persistente}: ao recuperar, o TM
lê o log e conclui as transações pendentes, chamando \texttt{xa\_recover()} em cada RM para obter
a lista de branches preparados. O mapeamento das interfaces TX e XA é direto:
\texttt{tx\_begin()} desencadeia \texttt{xa\_start()}; \texttt{tx\_commit()} desencadeia
\texttt{xa\_end()} + \texttt{xa\_prepare()} + \texttt{xa\_commit()} (ou
\texttt{xa\_rollback()} em caso de veto); \texttt{tx\_rollback()} desencadeia
\texttt{xa\_end()} + \texttt{xa\_rollback()}.

A coordenação da concorrência é assegurada pelo protocolo \textit{Two-Phase Locking} (2PL), que
garante serializabilidade dos escalonamentos concorrentes. O 2PL define duas fases: na
\textbf{fase de crescimento} (\textit{growing phase}), uma transação adquire locks mas nunca os
liberta; na \textbf{fase de contração} (\textit{shrinking phase}), liberta locks mas nunca adquire
novos. O ponto de transição — o \textit{lock point} — coincide com o commit ou rollback.

O sistema distingue dois tipos de lock: \textbf{ReadLock} (partilhado, S) para operações de
leitura, e \textbf{WriteLock} (exclusivo, X) para operações de escrita. Dois ReadLocks são
compatíveis entre transações concorrentes; qualquer combinação envolvendo um WriteLock gera
conflito e bloqueia a transação solicitante até a detentora libertar o recurso. A granularidade
dos locks é ao nível do par \texttt{(vecId,~index)} — cada posição de cada vetor é um recurso
independente — minimizando a contenção entre transações que acedem a elementos distintos.

A verificação do invariante pelo \texttt{Service.Validator} ilustra concretamente a necessidade do
2PL: para calcular a soma total de todos os elementos, o Validator inicia uma transação, adquire
ReadLocks em \emph{todos} os pares \texttt{(vecId, index)} e só então lê os valores e calcula
o somatório. Sem esta proteção, uma transferência em execução poderia ter já debitado o valor na
origem mas ainda não creditado no destino, produzindo uma leitura inconsistente. Com 2PL, os
ReadLocks bloqueiam qualquer WriteLock concorrente nesses elementos, garantindo que a leitura é
serializada. A deteção de deadlocks é realizada por timeout configurável: uma transação que não
consiga adquirir os locks necessários dentro do intervalo definido é abortada, libertando os seus
recursos e permitindo o progresso das transações concorrentes.

% -------------------------------------------------------
\subsection*{2.2 Dificuldade de Desenvolvimento com Múltiplos Elementos Serviço}
\addcontentsline{toc}{subsection}{2.2 Dificuldade de Desenvolvimento com Múltiplos Elementos Serviço}

O desenvolvimento de sistemas informáticos compostos por múltiplos elementos Serviço independentes
apresenta desafios que crescem super-linearmente com o número de componentes. Ao contrário de
um sistema monolítico, em que todas as unidades de execução partilham o mesmo espaço de
endereçamento e as falhas tendem a ser totais, num sistema distribuído as \textbf{falhas são
parciais}: um serviço pode estar ativo enquanto outro se torna inacessível, e não existe forma de
distinguir um nó lento de um nó falhado sem o recurso a timeouts — a chamada
\textit{impossibilidade de deteção de falhas} em modelos assíncronos.

No modelo de programação por \textbf{invocação remota} (RPC/gRPC), a sintaxe das chamadas
aproxima-se das chamadas locais, mas a semântica é fundamentalmente diferente: latência variável,
risco de duplicação de mensagens (\textit{at-least-once}), e necessidade de tratamento explícito
de timeouts e retentativas. Cada chamada gRPC entre serviços é um ponto de falha potencial que
exige decisão programática: abortar, retentar, ou propagar o erro.

A \textbf{gestão de estado distribuído} acrescenta complexidade adicional. Quando o estado
relevante para uma operação está disperso por $M$ instâncias de \texttt{Service.Vector}, qualquer
operação que abranja mais do que uma instância deve coordenar-se com o TM e o TPLM para garantir
consistência. O número de mensagens de coordenação por transação cresce com o número de RMs
participantes, e a probabilidade de falha durante o protocolo aumenta proporcionalmente.

A \textbf{heterogeneidade} dos elementos exacerba estas dificuldades: diferentes implementações
podem ter políticas distintas de gestão de threads, de timeouts, e de comportamento perante erros
de rede. Garantir que um sistema composto por $N$ serviços desenvolvidos por equipas distintas se
comporta corretamente perante falhas parciais requer uma disciplina rigorosa de contratos de
interface — papel desempenhado, neste trabalho, pelas especificações TX e XA e pelos ficheiros
\texttt{.proto} do gRPC. À medida que o número de serviços e as suas interdependências crescem,
a complexidade de raciocinar sobre o comportamento global do sistema em presença de falhas atinge
facilmente proporções intratáveis sem uma arquitetura de coordenação bem definida.

% -------------------------------------------------------
\subsection*{2.3 Aplicabilidade de MOM na Coordenação}
\addcontentsline{toc}{subsection}{2.3 Aplicabilidade de MOM na Coordenação}

A adoção de um \textit{Message Oriented Middleware} (MOM) --- como Apache Kafka ou RabbitMQ ---
na interação entre \texttt{Service.Client}, \texttt{Service.Vector}, \texttt{Service.TM} e
\texttt{Service.TPLM} apresenta vantagens relevantes para certos padrões de comunicação. O
\textbf{desacoplamento temporal} permite que o produtor e o consumidor de mensagens não estejam
disponíveis em simultâneo: um cliente poderia enfileirar pedidos de transferência mesmo que o TM
estivesse momentaneamente indisponível. O \textbf{desacoplamento espacial} liberta o produtor do
conhecimento do endereço do consumidor, e o \textbf{buffering} de mensagens absorve picos de
carga sem perda de pedidos.

No entanto, a utilização de MOM no contexto da coordenação transacional com 2PC apresenta
incompatibilidades fundamentais. O protocolo 2PC requer \textbf{confirmação explícita e ordenada
de cada fase}: o TM não pode emitir \texttt{xa\_commit()} sem ter recebido previamente resposta
de \texttt{xa\_prepare()} de todos os RMs. A semântica \textit{at-least-once} típica de um MOM
--- em que uma mensagem pode ser entregue mais de uma vez em caso de falha --- é incompatível
com a natureza das operações XA, a menos que se implemente deduplicação ao nível da aplicação,
o que acrescenta complexidade considerável. A semântica \textit{exactly-once}, disponível em
condições restritas em sistemas como o Kafka, introduz latência adicional e dependência de
funcionalidades de transação internas ao MOM, criando um segundo nível de coordenação transacional.

Por estas razões, este trabalho adota \textbf{comunicação gRPC síncrona} como protocolo de
coordenação: a natureza \textit{request-response} mapeia naturalmente para o fluxo de
confirmações do 2PC, e a definição de contratos em Protobuf garante serializações eficientes e
determinísticas. O MOM poderia ser equacionado para fluxos de menor criticidade transacional ---
por exemplo, notificações assíncronas de estado após a conclusão de uma transação --- mas não
como transporte primário da coordenação 2PC.

% -------------------------------------------------------
\subsection*{2.4 Heterogeneidade Tecnológica entre Elementos Serviço}
\addcontentsline{toc}{subsection}{2.4 Heterogeneidade Tecnológica entre Elementos Serviço}

Um requisito implícito de qualquer sistema distribuído aberto é a possibilidade de os seus
elementos serem desenvolvidos por equipas com culturas tecnológicas distintas. Num cenário
empresarial real, o \texttt{Service.Vector} poderia ser implementado em Java/Quarkus por uma
equipa, o \texttt{Service.TM} em Python por outra, e o \texttt{Service.Client} em C\#/.NET por
uma terceira, refletindo escolhas motivadas por competências internas, ecossistemas de bibliotecas
disponíveis, ou requisitos de desempenho específicos.

O protocolo \textbf{gRPC} viabiliza esta interoperabilidade através da definição dos contratos de
interface em ficheiros \texttt{.proto} independentes da linguagem de implementação. O compilador
\texttt{protoc} gera stubs cliente e servidor para Java, Python, C\#, Go e outras linguagens,
todos partilhando a mesma serialização binária (Protocol Buffers). Um cliente Java pode assim
invocar um serviço Python sem qualquer adaptação adicional, desde que ambos respeitem o mesmo
contrato \texttt{.proto}. Esta abordagem é conceitualmente análoga ao princípio de
substituibilidade definido pela especificação XA: a estrutura \texttt{xa\_switch\_t} é o contrato
que qualquer RM deve publicar, independentemente da sua tecnologia interna, permitindo ao TM
trabalhar com RMs de fornecedores distintos sem recompilação. Os ficheiros \texttt{.proto}
desempenham esse mesmo papel de contrato standard neste sistema.

Os limites desta interoperabilidade merecem ser considerados: diferenças na gestão de threads
entre linguagens podem afetar o comportamento sob carga (o modelo de I/O assíncrono do Python
\textit{asyncio} vs.\ o modelo de threads do Java); a gestão de memória e o \textit{garbage
collection} podem introduzir latências não determinísticas; e as bibliotecas gRPC de diferentes
plataformas podem apresentar comportamentos distintos perante condições de erro de rede, exigindo
estratégias de retry e timeout cuidadosamente alinhadas entre todos os participantes.

% -------------------------------------------------------
\subsection*{2.5 Escalabilidade do Sistema}
\addcontentsline{toc}{subsection}{2.5 Escalabilidade do Sistema}

Considere-se um cenário em que são instanciados $K$ serviços cliente e $M$ serviços vetor, com
$K$ e $M$ de ordem elevada. Neste contexto, a arquitetura apresenta dois \textit{bottlenecks}
de escalabilidade estruturais.

O \textbf{\texttt{Service.TPLM}} é um coordenador centralizado de locks: com $K$ clientes a
solicitar locks concorrentemente, o número de pedidos por unidade de tempo cresce com $K$, e a
latência de cada pedido aumenta com a profundidade da fila de espera. Em cenários de alta
contenção --- transações que acedem ao mesmo par \texttt{(vecId, index)} --- o TPLM torna-se o
gargalo dominante do sistema. A escalabilidade horizontal exige \textbf{particionamento do espaço
de locks}: diferentes instâncias do TPLM seriam responsáveis por subconjuntos disjuntos de
recursos (por exemplo, por \texttt{vecId}), eliminando a contenção entre transações que acedem a
vetores distintos. Esta estratégia requer um mecanismo de encaminhamento para que o
\texttt{Service.Client} saiba a qual instância de TPLM dirigir cada pedido de lock.

O \textbf{\texttt{Service.TM}} é \textit{stateful} e centralizado: mantém o estado de todas as
transações ativas e coordena o 2PC com todos os RMs participantes. Num modelo ativo-passivo com
eleição de líder via ZooKeeper, a disponibilidade é garantida após falha, mas o throughput é
limitado pela capacidade de uma única instância ativa.

O \textbf{Apache ZooKeeper} desempenha um papel duplo neste contexto: fornece a infraestrutura de
\textbf{eleição de líder} que garante que o sistema continua operacional após falha de uma
instância do TPLM ou TM (usando z-nodes efémeros --- o nó que detém o z-node \texttt{/tplm/leader}
é o líder; quando falha, o z-node desaparece e os candidatos em espera iniciam nova eleição); e
pode ser usado como \textbf{registo de configuração} para as instâncias dos serviços, facilitando
a adição dinâmica de novos nós ao sistema em crescimento. O protocolo ZAB (\textit{ZooKeeper
Atomic Broadcast}) que sustenta o ZooKeeper é estruturalmente análogo ao 2PC do modelo X/Open:
o Leader propõe (\texttt{PROPOSAL}), os Followers aceitam (\texttt{ACK}), e o Leader confirma
(\texttt{COMMIT}) após obter ACKs de um quórum --- garantindo consistência mesmo em presença de
falhas parciais.

% -------------------------------------------------------
\subsection*{2.6 Transações ACID vs.\ LLT/Saga}
\addcontentsline{toc}{subsection}{2.6 Transações ACID vs.\ LLT/Saga}

As quatro propriedades ACID constituem o contrato de correção das transações no modelo X/Open
DTP. A \textbf{Atomicidade} é garantida pelo protocolo 2PC: ou todos os RMs fazem commit, ou
todos fazem rollback --- a presunção de rollback garante este comportamento mesmo em caso de
falha do TM. A \textbf{Consistência} é responsabilidade conjunta do AP (que não deve violar os
invariantes de domínio ao compor as operações) e do TPLM (que garante que não existem estados
intermédios visíveis a outras transações durante a execução). O \textbf{Isolamento} é garantido
pelo 2PL: nenhuma transação vê os efeitos parciais de outra, por força dos locks de leitura e
escrita detidos até ao commit. A \textbf{Durabilidade} exige que cada RM persista as suas
alterações em armazenamento estável antes de responder afirmativamente ao \texttt{xa\_prepare()},
garantindo que os efeitos de uma transação confirmada sobrevivem a falhas subsequentes do sistema.

Este modelo é adequado para transações de curta duração, como as transferências entre vetores
deste trabalho, em que os locks são detidos por milissegundos a segundos e o custo do 2PC é
amortizado pela brevidade da operação.

Em contrapartida, em cenários de longa duração --- por exemplo, a reserva coordenada de voo,
hotel e transfer numa agência de viagens, cujas confirmações externas podem demorar horas --- o
modelo ACID com 2PL apresenta problemas sérios: os locks mantidos durante períodos prolongados
bloqueiam todo o sistema, e um rollback global pode ser semanticamente impossível (um voo já
confirmado com o operador aéreo requer um processo de cancelamento, não uma simples reversão).

O padrão \textbf{Saga} (Garcia-Molina, 1987) responde a esta limitação decompondo uma
\textit{Long-Lived Transaction} (LLT) numa sequência de transações curtas e atómicas, cada
uma emparelhada com uma transação de \textbf{compensação} capaz de desfazer o seu efeito. Se
uma etapa intermédia falhar, as compensações das etapas anteriores são executadas em ordem
inversa. O preço é a \textbf{ausência de isolamento forte}: entre o commit de uma etapa e a
eventual execução da sua compensação, os efeitos intermédios são visíveis a outras transações,
podendo ocorrer anomalias de leitura não reproduzíveis no modelo ACID clássico. No contexto
deste trabalho, o modelo ACID é a escolha natural: as transferências são curtas, o invariante
de somatório constante exige isolamento forte, e não existem confirmações externas ao sistema
que tornem o rollback semanticamente impossível.

% -------------------------------------------------------
\subsection*{2.7 Service Registry e Transparência à Localização (Consul)}
\addcontentsline{toc}{subsection}{2.7 Service Registry e Transparência à Localização (Consul)}

A \textbf{transparência à localização} é um requisito crítico em sistemas distribuídos dinâmicos:
os clientes não devem precisar de conhecer os endereços físicos dos serviços que invocam,
especialmente quando esses serviços podem ser replicados, migrados ou substituídos. O projeto
\textbf{Consul}, desenvolvido pela HashiCorp, é um \textit{service registry} que responde a este
requisito: cada serviço regista-se no arranque com o seu endereço e porta, e os clientes
consultam o Consul para descobrir as instâncias disponíveis. O Consul oferece ainda
\textit{health checks} periódicos --- removendo automaticamente do registo instâncias que deixaram
de responder --- e a possibilidade de armazenar configuração distribuída acessível a todos os
serviços em tempo de execução.

No exemplo entregue em aula prática, os serviços registam-se no Consul via
\texttt{ConsulRegistrar} no arranque, e os clientes resolvem os endereços dinamicamente em cada
invocação. Esta abordagem simplifica a gestão do ciclo de vida dos serviços: adicionar ou remover
instâncias não requer reconfiguração dos clientes.

Neste trabalho, optou-se por \textbf{não implementar o Consul}: a configuração estática dos
endereços dos serviços via variáveis de ambiente é suficiente para o demonstrador, cujo número
de instâncias é fixo e controlado. Esta opção simplifica o \textit{setup} e elimina a dependência
de uma infraestrutura adicional, reduzindo o número de pontos de falha no contexto do
demonstrador.

O impacto desta opção na escalabilidade e na gestão do ciclo de vida é, no entanto, relevante
para um contexto de produção: com $K$ clientes e $M$ vetores dinâmicos, a configuração estática
tornaria o sistema operacionalmente inflexível --- qualquer alteração ao número de instâncias
exigiria reconfiguração e reinicialização dos clientes. O Consul (ou uma solução equivalente,
como o próprio ZooKeeper usado como registo de serviços) seria a abordagem natural para gerir o
ciclo de vida em produção, assegurando que novos nós são descobertos automaticamente e que
instâncias falhadas são removidas sem intervenção manual.
